Thermal runaway reactions are represented using reduced-order kinetic
models that relate reaction progress to temperature and reaction state.

For reaction $i$, the general progress equation is

\begin{equation}
    \frac{d\alpha_i}{dt}
    =
    k_i(T)
    f_i
    \left(
        \alpha_i,
        \mathbf{\alpha}
    \right),
    \label{eq:general_kinetics}
\end{equation}

where $\alpha_i$ is the dimensionless reaction conversion,
$k_i(T)$ is the intrinsic temperature-dependent rate coefficient, and
$f_i$ is a progress function that may depend on the state of the
complete reaction network.


\subsection{Reaction conversion}

Reaction progress is represented by a conversion variable satisfying

\begin{equation}
    0
    \leq
    \alpha
    \leq
    1.
\end{equation}

The unreacted or remaining fraction is therefore

\begin{equation}
    r
    =
    1-\alpha.
\end{equation}

An unreacted state corresponds to $\alpha=0$, while a fully reacted
state corresponds to $\alpha=1$.


\subsection{Arrhenius kinetics}

The principal kinetic model currently implemented in LiionTR is the
Arrhenius relation,

\begin{equation}
    k(T)
    =
    A
    \exp
    \left(
        -\frac{E_a}{RT}
    \right),
    \label{eq:arrhenius}
\end{equation}

where $A$ is the pre-exponential factor, $E_a$ is the activation energy,
$R$ is the universal gas constant, and $T$ is the absolute temperature.

For first-order kinetics, $A$ has units of \si{\per\second}, while
$E_a$ is expressed in \si{\joule\per\mole}.

The exponential temperature dependence is central to thermal runaway.
A relatively modest increase in temperature can produce a large increase
in reaction rate, thereby increasing heat generation and further
accelerating the thermal response.


\subsection{Power-law progress}

A common progress model is based on the remaining reactant fraction,

\begin{equation}
    f(\alpha)
    =
    \left(
        1-\alpha
    \right)^n,
    \label{eq:power_law_progress}
\end{equation}

where $n$ is the reaction order.

The complete progress equation becomes

\begin{equation}
    \frac{d\alpha}{dt}
    =
    A
    \exp
    \left(
        -\frac{E_a}{RT}
    \right)
    \left(
        1-\alpha
    \right)^n.
\end{equation}

For $n=1$, the reaction rate decreases linearly with the remaining
reactant fraction.


\subsection{Autocatalytic progress}

LiionTR also supports autocatalytic progress laws of the general form

\begin{equation}
    f(\alpha)
    =
    \alpha^m
    \left(
        1-\alpha
    \right)^n,
    \label{eq:autocatalytic_progress}
\end{equation}

where $m$ and $n$ control the dependence on reacted and unreacted
material, respectively.

Such a formulation can represent reactions that accelerate as reaction
products or active sites accumulate before eventually slowing as the
reactant is depleted.

A frequently used special case is

\begin{equation}
    f(\alpha)
    =
    \alpha
    \left(
        1-\alpha
    \right).
\end{equation}


\subsection{Temperature thresholds}

Some reduced-order models assume that a reaction is inactive below a
prescribed onset temperature.

LiionTR represents this using a thresholded kinetic rate,

\begin{equation}
    k_{\mathrm{th}}(T)
    =
    \begin{cases}
        0,
        & T < T_{\mathrm{th}},
        \\[6pt]
        k(T),
        & T \geq T_{\mathrm{th}}.
    \end{cases}
    \label{eq:temperature_threshold}
\end{equation}

This formulation is useful when reproducing empirical literature models
that prescribe an explicit reaction onset temperature.


\subsection{Context-dependent reaction progress}

Thermal runaway reactions are not always independent.

The progress of one reaction may depend on the state of another
reaction, such as the depletion of a protective layer or the formation
of reactive products.

LiionTR therefore allows a progress model to depend on a reaction
context containing the complete conversion vector,

\begin{equation}
    \boldsymbol{\alpha}
    =
    \left[
        \alpha_1,
        \alpha_2,
        \ldots,
        \alpha_N
    \right]^{\mathrm{T}}.
\end{equation}

The progress equation can therefore be written more generally as

\begin{equation}
    \frac{d\alpha_i}{dt}
    =
    k_i(T)
    f_i
    \left(
        \alpha_i,
        \boldsymbol{\alpha}
    \right).
\end{equation}


\subsection{Threshold progress}

A context-dependent reaction may be activated only when another state
variable satisfies a prescribed condition.

For a context variable $x$, a generic threshold progress model may be
written as

\begin{equation}
    f_{\mathrm{th}}
    =
    \begin{cases}
        0,
        & x > x_{\mathrm{th}},
        \\[6pt]
        f(\alpha),
        & x \leq x_{\mathrm{th}}.
    \end{cases}
    \label{eq:progress_threshold}
\end{equation}

This formulation is used to represent sequential or conditionally
activated thermal runaway mechanisms.


\subsection{Exponential inhibition}

LiionTR also provides an exponential inhibition mechanism,

\begin{equation}
    f_{\mathrm{inh}}
    =
    f(\alpha)
    \exp(-x),
    \label{eq:exponential_inhibition}
\end{equation}

where $x$ is a context-dependent inhibition variable.

As $x$ increases, the effective reaction rate decreases exponentially.

Such a formulation can represent the influence of a protective layer or
other inhibiting material on a subsequent reaction.


\subsection{Normalized context variables}

A context variable may be constructed from the conversion or remaining
fraction of another reaction.

For example, the remaining fraction of reaction $j$ is

\begin{equation}
    r_j
    =
    1-\alpha_j.
\end{equation}

A normalized remaining-fraction variable may then be defined as

\begin{equation}
    x_j
    =
    \frac{
        r_j
    }{
        r_{\mathrm{ref}}
    },
    \label{eq:normalized_remaining_fraction}
\end{equation}

where $r_{\mathrm{ref}}$ is a prescribed reference value.

This provides a flexible mechanism for reproducing coupling terms
reported in reduced-order thermal runaway models.


\subsection{Separation of kinetics and progress}

LiionTR deliberately separates the intrinsic kinetic law from the
reaction-progress model.

The kinetic model determines how temperature affects the intrinsic rate,

\begin{equation}
    T
    \longrightarrow
    k(T),
\end{equation}

while the progress model determines how reaction state modifies that
rate,

\begin{equation}
    \boldsymbol{\alpha}
    \longrightarrow
    f
    \left(
        \alpha,
        \boldsymbol{\alpha}
    \right).
\end{equation}

The complete progress rate is obtained only after combining both terms.

This separation permits a common kinetic model, such as Arrhenius
kinetics, to be reused with multiple progress laws.


\subsection{Current limitations}

The present kinetic framework is intended primarily for reduced-order
thermal runaway modeling.

Current limitations include:

\begin{itemize}
    \item no explicit pressure dependence in the standard kinetic laws;
    \item no transport-limited reaction formulation;
    \item no explicit dependence on electrode potential or lithium
          concentration;
    \item no uncertainty propagation for kinetic parameters;
    \item no automatic kinetic parameter estimation from calorimetric
          data;
    \item no automatic competition between alternative reaction
          mechanisms.
\end{itemize}

Future developments may include parameter calibration using ARC or DSC
data, uncertainty quantification, and coupling to electrochemical state
variables supplied by PyBaMM.